\section*{Chapitre \ref{ch:s2:systematic-macro} --- Macro systématique}

\begin{solution}{exo:s2:systematic-macro:1}
$(0.84 + 0.78 + 0.33 + 0.29)/\sqrt4 = 1.12$.
\end{solution}

\begin{solution}{exo:s2:systematic-macro:2}
$(0.84 + 0.78)/\sqrt{2 + 2 \times (-0.43)} = 1.52$~: une corrélation négative fait monter le ratio de Sharpe de la combinaison au-dessus de la somme divisée par $\sqrt2$.
\end{solution}

\begin{solution}{exo:s2:systematic-macro:3}
$1 - 0.5^{1/3} = 21~\%$.
\end{solution}

\begin{solution}{exo:s2:systematic-macro:4}
La tendance dure environ un an, et le mouvement de prix sur cinq ans en contient encore la partie récente~: le retrancher revient à être vendeur sur les marchés qui ont suivi une tendance haussière,
à l'opposé du momentum. Le signal fondé sur l'ancrage ne comporte pas un tel recouvrement.
\end{solution}

\begin{solution}{exo:s2:systematic-macro:5}
Les données macroéconomiques sont révisées, parfois fortement~; un backtest sur données révisées utilise des chiffres que personne ne connaissait à l'époque et surestime le signal.
\end{solution}

\begin{solution}{exo:s2:systematic-macro:6}
Les classes d'actifs ont des volatilités et des facteurs communs différents~; au sein d'une classe, le portefeuille est neutre à la direction propre de la classe et ne négocie que le
signal, et le risque de chaque classe peut faire l'objet d'un budget distinct.
\end{solution}

\begin{solution}{exo:s2:systematic-macro:7}
La value tombe de 0.84 à 0.69 et la combinaison de 1.15 à 1.05~: un ancrage plus bruité dilue la value, et les trois autres signaux portent le portefeuille.
\end{solution}

\begin{solution}{exo:s2:systematic-macro:8}
Le ratio de Sharpe provient d'un marché synthétique et de signaux dont les corrélations peuvent monter ensemble en crise~; à 30~\% de volatilité, le pire
drawdown de la combinaison, de 21~\% à 10~\% de volatilité, serait d'environ 63~\%. Dimensionner les positions sur les corrélations et les drawdowns en situation de stress, non sur le ratio mesuré.
\end{solution}

\begin{solution}{pb:s2:systematic-macro:1}
\begin{enumerate}
\item Prix rapporté à un ancrage fondamental~: parité de pouvoir d'achat, rendements réels, rendements des bénéfices.
\item Position longue sur les tendances macro en amélioration, courte sur celles en dégradation~; publications de données comparées aux prévisions.
\item Long-short par rang au sein de chaque classe, mensuel, ajusté à 10~\% de volatilité.
\item Primes de value et de momentum partout, corrélées négativement, avec des facteurs communs aux classes d'actifs.
\item Un écart de 10~\% d'écart-type qui se résorbe en trois ans~; un indice de surprises statistiques qui se résorbe en 60 jours.
\item L'écart mesuré par rapport à un ancrage affecté d'une erreur persistante de 10~\%~; la surprise affectée d'un bruit de moitié de sa taille.
\item Il dilue le signal~: à 30~\% d'erreur, la value atteint 0.69.
\item \omterm{def:s2:systematic-macro:momentum}{Momentum macro}~: position longue sur les fondamentaux en amélioration, courte sur ceux qui se dégradent, avec une faible corrélation aux classes d'actifs.
\item Value 0.84, retournement $-0.84$, momentum 0.78, carry 0.33, surprise 0.29.
\item Proche de zéro sauf retournement et momentum, $-0.43$.
\item 1.15 avec la value fondée sur l'ancrage, 0.31 avec la value fondée sur le prix, 0.81 sans value.
\item La value fondée sur les prix recoupe la tendance et en devient le miroir.
\item \textbf{Résultat nommé.} Value 0.84, momentum 0.78, carry 0.33 et surprise 0.29, presque non corrélés, se combinent en 1.15 avec un drawdown plus faible
que chacun pris isolément~; la value fondée sur les prix est le miroir du momentum (corrélation $-0.43$) et tire la combinaison vers 0.31.
\item La somme des ratios de Sharpe divisée par la racine carrée de leur nombre~: 1.12.
\item Les corrélations ont tendance à monter~; le bénéfice de la combinaison se réduit justement quand on en a le plus besoin.
\item Pondération égale au départ, puis fondée sur les corrélations et volatilités glissantes, plafonnée.
\item Données point-in-time, signaux fixés à l'avance, coûts, plusieurs graines aléatoires ou plusieurs périodes.
\item Le biais en faveur du \omterm{def:s2:systematic-macro:surprise}{signal de surprises économiques}.
\item Il ajoute la value et les données macro à la tendance et au carry du livre~8.
\item Chacune est rémunérée pour un risque différent, et ces risques ne surviennent pas ensemble.
\end{enumerate}
\end{solution}

\begin{solution}{iq:s2:systematic-macro:1}
Devises~: taux de change réels face à la parité de pouvoir d'achat~; obligations~: rendements réels ou primes de terme par rapport à leur historique~; indices actions~: rendements des bénéfices ou rendements
comptables face aux rendements obligataires ou à leur propre historique.
\end{solution}

\begin{solution}{iq:s2:systematic-macro:2}
Le momentum achète ce qui a monté, la value ce qui a baissé par rapport à son ancrage, si bien qu'ils prennent souvent des positions opposées~; la corrélation négative rend la
combinaison plus stable que chacun des deux.
\end{solution}

\begin{solution}{iq:s2:systematic-macro:3}
Il faut les compenser selon les budgets de risque du portefeuille~; les signaux sont conçus pour parfois se contredire, et passer outre l'un parce que l'autre existe rompt la
diversification de la combinaison.
\end{solution}

\begin{solution}{iq:s2:systematic-macro:4}
Un retournement de tendance sur tous les marchés, un krach du carry, un pic de corrélations entre classes d'actifs, et un choc de données (une grande surprise statistique) qui fait bouger tous
les pays ensemble.
\end{solution}

\begin{solution}{iq:s2:systematic-macro:5}
Conserver chaque publication avec sa date et son heure, sa première version et chacune de ses révisions, ainsi que la prévision consensuelle qui la précède~; interroger selon la date de connaissance, non
la période de référence.
\end{solution}

\begin{solution}{iq:s2:systematic-macro:6}
Chaque signal $i$ est ajusté à la volatilité $\sigma$ et rapporte $S_i\sigma$~; la somme à pondération égale rapporte $\sigma\sum S_i$ avec une variance $\sigma^2(n +
n(n - 1)\rho)$, soit un ratio égal à $\sum_i S_i/\sqrt{n + n(n - 1)\rho}$.
\end{solution}
